\chapter{CONCLUSION AND FUTURE WORK}
\label{chap:conclusion}

\section{Summary of contributions}

This thesis designed and implemented an Amazon Marketplace Intelligence Tracker for scheduled competitive monitoring of consumer-electronics products. The contribution is primarily an integration and systems contribution: it connects Amazon data extraction, relational storage, analytical scripts, dashboard exploration, and Telegram-based agent delivery into one reproducible workflow.

The implemented framework addresses the gap between raw marketplace data and usable competitive intelligence. Apify actors collect category rankings, watchlist product snapshots, and customer reviews. Supabase stores both raw and derived tables. Python scripts compute entrant/exit events, sponsored share of voice, price tiers, sentiment, listing quality, brand momentum, alerts, image-change flags, and price forecasts. The Streamlit and web dashboards support pull-based analysis, while three OpenClaw Telegram agents provide push-based summaries through fourteen read-only skills.

The most important design choice is the separation between data mutation and conversational delivery. Ingestion and analytics scripts write to the database on schedule, while OpenClaw skills only read existing Supabase records. This makes the agent layer useful for communication without allowing it to alter the analytical state of the system. For a thesis-scale deployment, this architecture is easier to audit and more defensible than allowing LLM agents to directly mutate the warehouse.

\section{Key empirical findings}

The 13-day pilot collected 1,756 category-ranking rows, 409 daily product snapshots, and 3,986 raw review records. These data supported both descriptive market monitoring and lightweight model evaluation.

For sentiment analysis, RoBERTa achieved 81.17\% accuracy and a macro-$F_1$ score of 0.594 on 3,851 reviews evaluated with star ratings as a distant-supervision proxy. The model performed well on positive reviews ($F_1 = 0.908$) and achieved useful negative-review recall (0.782), which is valuable for detecting customer complaints. However, the neutral class remained weak ($F_1 = 0.158$), confirming that 3-star reviews are a noisy proxy for neutral textual sentiment.

For price forecasting, a five-way temporal holdout comparison was conducted across all 30 watchlist ASINs with a three-day holdout. The models evaluated were: naive carry-forward, linear trend (OLS), Exponential Smoothing (ETS), ARIMA(1,1,0), and Prophet. The naive baseline achieved the lowest pooled MAE (\$1.292), with ARIMA the closest non-trivial alternative (\$1.323, only 2.4\% higher). Prophet ranked fifth with MAE of \$2.970 and MAPE of 6.75\%. The near-equivalence of the naive and ARIMA results confirms that Amazon prices approximate a random walk over short horizons — a data-volume constraint rather than a model deficiency. Despite its lower point-forecast accuracy, Prophet is retained in the deployed system because it is the only evaluated model that produces native uncertainty intervals; its empirical 80\% prediction-interval coverage was 67.8\%, falling below the nominal target and indicating undercalibration on the current short training windows.

The rule-based analytics generated 1,060 entrant/exit events and 1,248 alert records, including 55 high-severity alerts. The alert engine helped separate routine Top-50 churn from more important events such as top-10 entrants, price drops, and stockout signals. The LQS results showed that many tracked listings were already highly optimized, with an average score of 93.4 out of 100, so LQS is most useful for finding weaker outliers rather than explaining all rank movement. The image-change module was implemented, but no observed image change exceeded the configured pHash threshold during the pilot.

\section{Limitations of the study}

Several limitations should be considered when interpreting the results.

\subsection{Short observation window}

The most important limitation is the 13-day observation window. This period is enough to demonstrate ingestion, storage, analytics, dashboarding, and agent delivery, but it is too short for robust time-series modelling. Weekly and yearly Prophet seasonality were intentionally disabled. The five-way forecast evaluation confirmed that no model outperformed the naive carry-forward baseline, with ARIMA(1,1,0) being the best non-trivial model (MAE \$1.323 versus naive \$1.292). This is consistent with the near-random-walk behavior of Amazon prices over short horizons. Longer data collection — ideally 60 or more days — would be required before the forecast component can be used as more than a short-term trend and volatility indicator.

\subsection{Noisy sentiment ground truth}

The primary sentiment evaluation on $N = 3{,}851$ reviews used star ratings as a distant-supervision proxy label. While 200 reviews were manually labelled during this project to quantify proxy noise (Table~\ref{tab:manual_label_eval}), this sample is too small to replace the proxy-based evaluation for the full corpus. The proxy introduces label noise, especially for 3-star reviews. A three-star review may contain both praise and criticism, or may penalize shipping rather than the product itself. The weak neutral-class result should therefore be interpreted as both a model limitation and a limitation of the proxy-labelling method.

\subsection{Review corpus sampling and sentiment signal staleness}

Two related limitations affect the reliability of the sentiment component. First, the review corpus used for RoBERTa inference is a sampled subset of available Amazon reviews rather than an exhaustive collection. The Apify Amazon Review Scraper actor retrieved approximately 3,986 review records across all 30 watchlist ASINs over the 13-day pilot period, yielding an average of roughly 133 reviews per ASIN. This sampling is a consequence of platform-level rate limiting and anti-scraping measures enforced by Amazon: the actor can retrieve only a limited number of review pages per run, so high-volume products with thousands of published reviews are represented by a small fraction of their actual review base. Sentiment scores produced by the system therefore reflect a convenience sample rather than the full distribution of customer opinion, and may over-represent recent or highly-ranked reviews depending on the ordering returned by the actor. This constraint should be disclosed whenever sentiment scores are cited as evidence of customer perception.

Second, there is an important architectural distinction between two review-related signals in the system. The Brand Momentum Score (BMS) component \texttt{review\_velocity} is computed as the day-over-day delta of \texttt{reviews\_count} in the \texttt{daily\_snapshots} table --- that is, the change in the total review count displayed on the Amazon product page, captured by the Amazon Product Scraper actor on each daily run. This signal is unaffected by review text scraping. By contrast, the sentiment score component of BMS and all outputs of the Sentiment Detective agent depend on the \texttt{reviews\_raw} table populated by the Amazon Review Scraper actor. During the pilot, review text ingestion ceased after 19~April 2026 due to a scheduling gap in the review scraping workflow. As a result, sentiment scores reported from that date onward are computed from historical reviews only and do not capture any shift in customer opinion that may have occurred during the final weeks of the observation window. The sentiment signal in the deployed brief should therefore be treated as a lagging indicator rather than a current measure.

\subsection{Deployment and engineering constraints}

The prototype was deployed as a thesis pilot rather than a hardened commercial service. Telegram polling can occasionally lag in the VMware environment, structured logging and formal unit tests are limited, and Supabase row-level security for the public web dashboard remains a configuration follow-up. Forecast data now persists to Supabase, but the system still depends on scheduled jobs finishing as expected. These constraints do not invalidate the prototype, but they define the boundary between a working academic system and a hardened service.

An attempt to validate LQS internal consistency via Cronbach's alpha was uninformative due to ceiling effects: three of seven components (title score, bullet score, and image score) exhibited zero variance across the curated watchlist, precluding meaningful computation. Similarly, BMS forward predictive validity (per-ASIN Spearman $\rho$, lag = 5 days) was not confirmed, consistent with its framing as a descriptive ranking heuristic rather than a predictive model.

\section{Future work}

Future work should focus on strengthening the parts of the system that were most constrained during the thesis timeline.

\subsection{Longer data collection and improved forecasting accuracy}

The first priority is to extend the data window from days to months. A longer history would allow the project to evaluate weekly cycles, promotional periods, and event-driven price changes more fairly. The thesis already conducted a five-way model comparison (naive, OLS, ETS, ARIMA, Prophet); the next step is to re-run this evaluation after accumulating 60 or more days of data, at which point trend-aware models such as ETS and Prophet are expected to gain a meaningful advantage over the naive baseline.

\subsection{Expanded manual sentiment annotation}

A 200-review manually labelled validation set was completed during this project to quantify proxy label noise (Table~\ref{tab:manual_label_eval}). The results confirmed that star-proxy labels carry approximately 25.5\% noise on this corpus, and that RoBERTa's true accuracy against manual labels (72.0\%) is 8.5 percentage points higher than its apparent proxy-evaluation accuracy. Expanding manual annotation to at least 1,000 reviews with a second independent annotator would enable inter-annotator agreement scoring and yield a dataset large enough to support fine-tuning or few-shot adaptation of the sentiment model. Aspect-level annotation --- distinguishing product-quality sentiment from complaints about shipping, delivery, or expectation mismatch --- would further increase the specificity of the deployed sentiment signals.

\subsection{Operational hardening}

The engineering layer can be strengthened through formal unit tests, structured logs, data-quality checks, and transactional alert upserts. Supabase row-level security should be tightened before exposing the public web dashboard broadly. The OpenClaw skills should remain read-only, but their outputs could be evaluated with a fixed set of benchmark questions to measure factual accuracy, brevity, and usefulness.

\subsection{Richer image and listing-change explanation}

The current pHash method can detect visual change but cannot explain what changed. A realistic next step is to combine pHash with a vision-language model or a rule-based image metadata comparison so that the system can distinguish between a new product angle, a promotional badge, a packaging redesign, or a minor crop. This should be added only after enough image-change events are collected to evaluate whether the added complexity is useful.

\subsection{Agent evaluation and workflow design}

The three-agent Telegram design can be expanded with systematic user evaluation. Future work could compare dashboard-only use, agent-only briefs, and combined dashboard-plus-agent workflows for decision speed and perceived usefulness. This would move the research contribution beyond technical feasibility toward evidence about how small sellers or analysts actually use the intelligence produced by the system.
